\Article{33}
Because of the importance of the matter, and so that everything may be ready for calculation, it seemed appropriate also to write down here explicit formulas for determining the coefficients \(A',A'',A'''\) etc., \(B'',B'''\) etc. etc. This calculation can be carried out in two ways, since identities must result both when the values of \(u^0,u',u''\) etc. taken from III are substituted in II, and when the values of \(\xi,\eta,\zeta\) etc. from II are substituted in III. The first way supplies the following systems of formulas:
\[
\begin{aligned}
\frac{\mathfrak{B}^0}{\mathfrak{A}^0}+A'&=0\\[4pt]
\frac{\mathfrak{C}^0}{\mathfrak{A}^0}+\frac{\mathfrak{C}'}{\mathfrak{B}'}.A'+A''&=0\\[4pt]
\frac{\mathfrak{D}^0}{\mathfrak{A}^0}+\frac{\mathfrak{D}'}{\mathfrak{B}'}.A'+\frac{\mathfrak{D}''}{\mathfrak{C}''}.A''+A'''&=0
\end{aligned}
\]
etc., from which \(A',A'',A'''\) etc. are found.
\[
\begin{aligned}
\frac{\mathfrak{C}'}{\mathfrak{B}'}+B''&=0\\[4pt]
\frac{\mathfrak{D}'}{\mathfrak{B}'}+\frac{\mathfrak{D}''}{\mathfrak{C}''}.B''+B'''&=0
\end{aligned}
\]
etc., from which \(B'',B'''\) etc. are found.
\[
\frac{\mathfrak{D}''}{\mathfrak{C}''}+C'''=0
\]
etc., from which \(C'''\) etc. are found. And so on.

The other way gives these formulas:
\[
\mathfrak{A}^0A'+\mathfrak{B}^0=0
\]
from which \(A'\) is obtained.
\[
\begin{array}{rllll}
\mathfrak{A}^0A''&+\mathfrak{B}^0B''&+\mathfrak{C}^0&=0\\
&\mathfrak{B}'B''&+\mathfrak{C}'&=0
\end{array}
\]
from which \(B''\) and \(A''\) are found.
\SourceLeaf{046}
\[
\begin{array}{rlllll}
\mathfrak{A}^0A'''&+\mathfrak{B}^0B'''&+\mathfrak{C}^0C'''&+\mathfrak{D}^0&=0\\
&\mathfrak{B}'B'''&+\mathfrak{C}'C'''&+\mathfrak{D}'&=0\\
&&\mathfrak{C}''C'''&+\mathfrak{D}''&=0
\end{array}
\]
from which \(C''',B''',A'''\) are found. And so on.

The two ways are almost equally convenient if the weights of the determinations of all the quantities \(x,y,z\) etc. are wanted; but whenever only one or two of the quantities \([\alpha\alpha],\allowbreak[\beta\beta],\allowbreak[\gamma\gamma]\) etc. are required, the first system will clearly be far preferable.

Moreover, combining equations I with IV leads to the same formulas and also supplies two ways of calculating the most plausible values \(A,B,C\) etc. themselves; namely, \textit{first},
\[
\begin{aligned}
A&=-\frac{\mathfrak{L}^0}{\mathfrak{A}^0}-A'\frac{\mathfrak{L}'}{\mathfrak{B}'}-A''\frac{\mathfrak{L}''}{\mathfrak{C}''}-A'''\frac{\mathfrak{L}'''}{\mathfrak{D}'''}-\text{etc.}\\[4pt]
B&=-\frac{\mathfrak{L}'}{\mathfrak{B}'}-B''\frac{\mathfrak{L}''}{\mathfrak{C}''}-B'''\frac{\mathfrak{L}'''}{\mathfrak{D}'''}-\text{etc.}\\[4pt]
C&=-\frac{\mathfrak{L}''}{\mathfrak{C}''}-C'''\frac{\mathfrak{L}'''}{\mathfrak{D}'''}-\text{etc.}\\
&\text{etc.}
\end{aligned}
\]
The other calculation is identical with the ordinary one, in which we put \(u^0=0,\allowbreak u'=0,\allowbreak u''=0\) etc.

\Article{34}
What we set out in art.\ 32 consists only of special cases of a more general theorem, which is as follows:

\textsc{Theorem.} \textit{Let \(t\) denote the following linear function of the indeterminates \(x,y,z\) etc.}
\[
t=fx+gy+hz+\text{etc.}+k
\]
\textit{which, when transformed into a function of the indeterminates \(u^0,u',u''\) etc., becomes}
\[
t=k^0u^0+k'u'+k''u''+\text{etc.}+K
\]
\textit{Then \(K\) will be the most plausible value of \(t\), and the weight of this determination will be}
\[
=\frac1{\mathfrak{A}^0k^0k^0+\mathfrak{B}'k'k'+\mathfrak{C}''k''k''+\text{etc.}}
\]
\textit{Proof.} The first part of the theorem is clear because the most plausible value of \(t\) must correspond to the values \(u^0=0,\allowbreak u'=0,\allowbreak u''=0\) etc. To prove the lat\SourceLeaf{047}%
ter part, we observe that, since \(\tfrac12d\Omega=\xi\,dx+\eta\,dy+\zeta\,dz+\text{etc.}\) and \(dt=f\,dx+g\,dy+h\,dz+\text{etc.}\), the following holds for \(\xi=f.\quad\eta=g,\allowbreak\zeta=h\) etc., independently of the values of the differentials \(dx,dy,dz\) etc.:
\[
d\Omega=2dt
\]
It therefore follows that, for the same values \(\xi=f,\allowbreak\eta=g,\allowbreak\zeta=h\) etc.,
\[
\begin{aligned}
&\frac{u^0}{\mathfrak{A}^0}du^0+\frac{u'}{\mathfrak{B}'}du'+\frac{u''}{\mathfrak{C}''}.du''+\text{etc.}\\
&\qquad=k^0du^0+k'du'+k''du''+\text{etc.}
\end{aligned}
\]
It is now readily seen that if \(dx,dy,dz\) etc. are independent of one another, then \(du^0,du',du''\) etc. are also independent of one another; from this we conclude that for \(\xi=f,\allowbreak\eta=g,\allowbreak\zeta=h\) etc.,
\[
u^0=\mathfrak{A}^0k^0,\quad u'=\mathfrak{B}'k',\quad u''=\mathfrak{C}''k''\quad\text{etc.}
\]
The value of \(\Omega\) corresponding to the same values will therefore be
\[
=\mathfrak{A}^0k^0k^0+\mathfrak{B}'k'k'+\mathfrak{C}''k''k''+\text{etc.}+M
\]
from which the truth of our theorem follows directly by art.\ 29.

If, moreover, we wish to carry out the transformation of the function \(t\) directly, i.\ e. without knowing substitutions IV of art.\ 32, the following formulas are available:
\[
\begin{aligned}
f&=\mathfrak{A}^0k^0\\
g&=\mathfrak{B}^0k^0+\mathfrak{B}'k'\\
h&=\mathfrak{C}^0k^0+\mathfrak{C}'k'+\mathfrak{C}''k''\quad\text{etc.,}
\end{aligned}
\]
from which the coefficients \(k^0,k',k''\) etc. will be determined successively, and finally we shall have
\[
K=k-\mathfrak{L}^0k^0-\mathfrak{L}'k'-\mathfrak{L}''k''-\text{etc.}
\]

\clearpage
